\section{Ordered pairs and products}
\label{sec:set-prerequisites-02}

\begin{definition}[Ordered pair]
The \emph{ordered pair} of sets \(a\) and \(b\) is
\[
  \tuple{a,b}\defeq\set{\set{a},\set{a,b}}.
\]
\end{definition}

\begin{lemma}[Ordered-pair identity]
For all sets \(a,b,c,d\),
\[
  \tuple{a,b}=\tuple{c,d}
  \quad\Longleftrightarrow\quad
  a=c\text{ and }b=d.
\]
\end{lemma}
\begin{proof}
The reverse implication follows by substitution.  For the forward implication,
if \(a=b\), then \(\tuple{a,b}=\set{\set a}\), so equality of the ordered
pairs forces \(c=d=a\).  If \(a\neq b\), then the ordered pair has the two
distinct members \(\set a\) and \(\set{a,b}\).  Equality forces
\(\set c=\set a\), hence \(c=a\), and then comparison of the two-element
members gives \(d=b\).  The remaining apparent case \(c=d\) would make the
right ordered pair a singleton and is excluded by \(a\neq b\).
\end{proof}

\begin{definition}[Cartesian product]
For sets \(A\) and \(B\),
\[
  A\times B\defeq\set{\tuple{a,b}:a\in A\text{ and }b\in B}.
\]
For every \(n\in\N\), define \(A^n\) as the set of functions from
\(n=\set{0,\ldots,n-1}\) to \(A\).  When \(n=0\), the domain is empty, so
\(A^0\) contains exactly the empty function.
\end{definition}

